% T24：只替换冻结版中两个局部问题；不修改正文文件。
% Appendix D 第128--132行，替换为：
在附录 B 定义的角动量投影与非正交度规下，可将归一化初、末态写为
\[
|\Psi_s^{I_sM_s}\rangle=\sum_{aK}f^{(s)}_{aK}P^{I_s}_{M_sK}|\Phi_a\rangle,
\quad (f^{(s)})^\dagger N^{I_s} f^{(s)}=1,\quad
N^I_{aK,bK'}=\langle\Phi_a|P^I_{KK'}|\Phi_b\rangle,\quad s=i,f.
\]
因此球张量的实验室系矩阵元是
\[
\langle\Psi_f^{I_fM_f}|T_{L\mu}|\Psi_i^{I_iM_i}\rangle
=\sum_{aK_f,bK_i}f^{(f)*}_{aK_f}f^{(i)}_{bK_i}
\langle\Phi_a|P^{I_f}_{K_fM_f}T_{L\mu}P^{I_i}_{M_iK_i}|\Phi_b\rangle.
\]
根据附录 A 式~\eqref{eq:appA-we-3j}，选取非零的 $3j$ 系数，即可恢复与磁量子数无关的约化矩阵元；若所有相应 $3j$ 系数均为零，则该角动量耦合不允许。范数核依照附录 B 的投影算符及重叠度规定义；式~\eqref{eq:appB-haar-measure} 仅定义 Haar 测度。投影可一致处理对称性恢复与非正交重叠，具体极限的有限性仍需对所用态和模型单独检查。

% Appendix D 第168行中“各张量分量通过……干涉”改为：
各张量分量与轻子波函数共同决定能量依赖的谱形状因子：
% 保留第169--171行既有公式。第172行第一句替换为：
对未极化初态、末态磁量子数未观测且轻子角度已积分的电子能谱，Zhi 等的展开将不同核总角动量秩的贡献按秩分块相加；同秩内不同算符矩阵元可以相干干涉，例如秩 0 的 $\zeta_0$ 与 $w$，以及秩 1 的 $x$ 与 $u$。此限定不排除极化或角关联观测量中的其他干涉。
% 第172行后续关于 N=126 的句子不在本次修订范围。

% sec-4-7 冻结版第71行，仅将“包含电子动能与不同秩矩阵元交叉干涉的复杂能量函数”替换为：
对未极化初态、末态磁量子数未观测且轻子角度已积分的电子能谱，依赖电子能量并由各核张量秩的贡献相加而成的函数；同秩内不同算符矩阵元可相干干涉
